%% Responsable: Sebastián López
\section*{Anexo: decisiones citadas}
\addcontentsline{toc}{section}{Anexo: decisiones citadas}
\label{sec:decisiones}

El texto cita entre paréntesis las decisiones del proyecto por su identificador (D1,
D2, \dots). Se tomaron antes de ejecutar y no cambian sin evidencia nueva. La tabla
resume las citadas; la lista completa, con su justificación, está en el repositorio
(\texttt{AGENTS.md}, sección 5).

\begin{center}
\small
\begin{tabular}{lp{0.8\textwidth}}
\toprule
ID & Decisión \\
\midrule
D1 & El repositorio base manda en datos y ejecución; el enunciado, en entregables, pesos y
rigor. Las diferencias se registran en la sección~\ref{sec:discrepancias}. \\
D2 & $L=5$, $k_i=8$, $|\conf|=32\,768$, $n=100$, validación con 87 instancias,
temperatura $0{,}0$, modelo \texttt{qwen17b} y semilla 1. \\
D3 & Técnicas elegidas: backtracking, búsqueda por coordenadas y reparto uniforme frente a
UCB1. \\
D4 & Successive Halving no se implementa: su costo supera el presupuesto y con pocas
instancias iniciales casi todas las configuraciones darían 0\,\%. \\
D5 & Cota del exacto: aciertos más instancias restantes; si no alcanza a la mejor, se
abandona la configuración. \\
D7 & El exacto se mide con el modelo sintético; el modelo real da los resultados de
precisión. \\
D9 & Reducción desde 3-SAT en su versión de decisión, con $L$ y $k_i$ como parte de la
entrada. \\
D11 & $B$ cuenta las generaciones nuevas por (configuración, instancia, semilla). \\
D12 & La validación con 87 instancias se usa una sola vez, sobre la configuración ya
elegida. \\
D13 & La guía de la primera entrega manda sobre el enunciado en esta entrega y en la
reducción. \\
D17 & Búsqueda real con $B=14\,400$ generaciones y dos arranques. \\
D18 & Tope de tres pasadas por arranque, el que se usó en los dos arranques (D17 preveía
dos). Se aprobó después de ejecutarlos. \\
\bottomrule
\end{tabular}
\end{center}
